% BHU PYQ -- Professional Exam Template
\documentclass[11pt,a4paper]{article}

\usepackage[a4paper,top=2.5cm,bottom=2.5cm,left=2.8cm,right=2.8cm,headheight=14pt]{geometry}
\usepackage{amsmath,amssymb}
\usepackage[T1]{fontenc}
\usepackage[utf8]{inputenc}
\usepackage{lmodern}
\usepackage{microtype}
\usepackage{fancyhdr}
\usepackage{enumitem}
\usepackage{hyperref}
\usepackage{lastpage}
\usepackage{parskip}
\usepackage{tikz}
\usetikzlibrary{arrows.meta,shapes,positioning}

\hypersetup{colorlinks=true,urlcolor=black,linkcolor=black,
  pdftitle={B.Sc. (Hons.) Zoology Sem-IV ZOB-403A -- BHU PYQ}}

\pagestyle{fancy}
\fancyhf{}
\renewcommand{\headrulewidth}{0.4pt}
\renewcommand{\footrulewidth}{0.4pt}
\fancyhead[L]{\small   \textit{Banaras Hindu University}}
\fancyhead[R]{\small   \textit{Zoology Paper: ZOB-403A}}
\fancyfoot[C]{\small Page  \thepage\ of \pageref{LastPage}}


\newlist{parts}{enumerate}{2}
\setlist[parts,1]{label=(\alph*),leftmargin=*,itemsep=4pt,topsep=3pt}
\setlist[parts,2]{label=(\roman*),leftmargin=*,itemsep=2pt,topsep=2pt}

\newcommand{\pts}[1]{\hfill   \textit{[#1]}}


\begin{document}


\begin{center}
  {\large\bfseries BANARAS HINDU UNIVERSITY}\\[2pt]
  {\normalsize B.Sc. (Hons.) Semester IV Examination, 2022-23}\\[6pt]
  \rule{0.8\linewidth}{0.5pt}\\[6pt]
  {\large\bfseries Zoology}\\[4pt]
  {\large\bfseries Paper: ZOB-403A --- Biology: Applied Zoology (Ancillary Course)}\\[4pt]
  {\normalsize Time: Three hours \hfill Full Marks: 70}
\end{center}

\rule{\linewidth}{0.4pt}

\smallskip



\noindent   \textit{Note: Answer six questions including Question No. 1 which is compulsory. The figures in the right-hand margin indicate marks.}

\bigskip

\medskip


\noindent   \textbf{Question 1.} Answer the following parts:

\begin{parts}
  \item State whether the following statements are True or False, giving reasons in 2-3 sentences each:\pts{$2  \times 6 = 12$}
  
\begin{parts}
    \item Lac is mentioned in   \textit{Mahabharata}.
    \item India doesn't offer a huge potential for aquaculture development.
    \item DDT is good for the environment.
    \item Spectacled Cobra is found in North-East India.
    \item   \textit{Cirrhinus cirrhosus} is among India's major indigenous carp.
    \item   \textit{Bombyx mori} produces Tasar silk.
  \end{parts}
  
  \item Define the following terms:\pts{$1  \times 6 = 6$}
  
\begin{parts}
    \item Beeswax
    \item Mariculture
    \item Polyculture
    \item Benthic effect
    \item Eutrophication
    \item Estuaries
  \end{parts}

  \item Differentiate between Krait and Cobra.\pts{2}
\end{parts}

\medskip


\noindent   \textbf{Question 2.} How is silk produced and what is its economic importance?\pts{10}

\medskip


\noindent   \textbf{Question 3.} Discuss the various tasks that worker honeybees perform in the colony. Discuss the way honey is prepared from nectar.\pts{10}

\medskip


\noindent   \textbf{Question 4.} How are transgenic animals made? Discuss their importance with one example.\pts{10}

\pagebreak
\medskip


\noindent   \textbf{Question 5.} Write a detailed note with illustration on the life cycle of the Lac insect.\pts{10}


\begin{center}

\begin{tikzpicture}[scale=1.0, >=Stealth, every node/.style={font=\small, align=center}]
  % Lifecycle of Lac Insect
  
ode (egg) at (0,3) [draw, rounded corners, fill=red!5] {   \textbf{Eggs}\\ (Laid in chambers)};
  
ode (nymph) at (3.5,1) [draw, rounded corners, fill=orange!5] {   \textbf{Crawlers (Nymphs)}\\ (Emergence \& search for host)};
  
ode (settle) at (3.5,-1.5) [draw, rounded corners, fill=orange!10] {   \textbf{Settlement}\\ (Attach to twigs)};
  
ode (lac) at (0,-3.5) [draw, rounded corners, fill=yellow!10] {   \textbf{Lac Secretion}\\ (Resin cell formation)};
  
ode (adults) at (-3.5,-1.5) [draw, rounded corners, fill=green!10] {   \textbf{Adult Differentiation}\\ (Wingless Female / Winged Male)};
  
ode (fertilization) at (-3.5,1) [draw, rounded corners, fill=blue!5] {   \textbf{Fertilization}\\ (Mating / Parthenogenesis)};

  % Flow arrows
  \draw[->, thick] (egg) -- (nymph);
  \draw[->, thick] (nymph) -- (settle);
  \draw[->, thick] (settle) -- (lac);
  \draw[->, thick] (lac) -- (adults);
  \draw[->, thick] (adults) -- (fertilization);
  \draw[->, thick] (fertilization) -- (egg);
\end{tikzpicture}
\end{center}

\medskip


\noindent   \textbf{Question 6.} Describe the reproductive system and the process of egg formation in a hen with the help of suitable diagrams.\pts{10}


\begin{center}

\begin{tikzpicture}[scale=0.9, >=Stealth, every node/.style={font=\small, align=center}]
  % Hen Oviduct Schematic
  
ode (ovary) at (-6.5,0) [draw, circle, fill=yellow!20] {   \textbf{Ovary}\\ (Releases yolk)};
  
ode (infundibulum) at (-3.5,0) [draw, rounded corners, fill=orange!5] {   \textbf{Infundibulum}\\ (Fertilization\\ 15 mins)};
  
ode (magnum) at (-0.5,0) [draw, rounded corners, fill=orange!10] {   \textbf{Magnum}\\ (Albumen added\\ 3 hours)};
  
ode (isthmus) at (2.5,0) [draw, rounded corners, fill=orange!15] {   \textbf{Isthmus}\\ (Shell membranes\\ 1.25 hours)};
  
ode (uterus) at (5.5,0) [draw, rounded corners, fill=red!10] {   \textbf{Uterus}\\ (Shell gland - Calcite\\ shell \& pigment\\ 20 hours)};
  
ode (vagina) at (8.5,0) [draw, rounded corners, fill=gray!15] {   \textbf{Vagina}\\ (Cuticle / Laying\\ 15 mins)};

  % Flows
  \draw[->, thick] (ovary) -- (infundibulum);
  \draw[->, thick] (infundibulum) -- (magnum);
  \draw[->, thick] (magnum) -- (isthmus);
  \draw[->, thick] (isthmus) -- (uterus);
  \draw[->, thick] (uterus) -- (vagina);
\end{tikzpicture}
\end{center}

\medskip


\noindent   \textbf{Question 7.} Write an essay on milk and milk products.\pts{10}

\medskip


\noindent   \textbf{Question 8.} Write short notes on the following:\pts{$2.5  \times 4 = 10$}

\begin{parts}
  \item Lac incrustation
  \item Brackish Water Aquaculture
  \item Attraction and repulsion cage mechanism
  \item Pearl formation
\end{parts}

\vfill
\rule{\linewidth}{0.4pt}

\begin{center}\small
\href{https://bhu-exam.vercel.app/}{Click Here}\\[4pt]\href{https://me-aryan.vercel.app/}{Click Here}\\[4pt]\href{https://quashan-me.vercel.app/}{Click Here}
\end{center}

\end{document}
